\section{Unique continuation theorems}\label{section-UCP}
As proved in Section~\ref{section-asymp-compact}, the last step to prove stabilization is the unique continuation property: if $z$ is a global solution of
\begin{equation}\label{eq-UCP2}
\begin{cases}
\partial_{tt}^2 z(x,t) = \Delta z(x,t) - \alpha z(x,t) - f'_u(x,w(x,t)) z\;&(x,t)\in\Omega\times\RR\\
z_{|\partial\Omega}(x,t)=0&(x,t)\in\partial\Omega\times \RR\\
\mkern20mu z(x,t)\equiv 0 & (x,t)\in \text{support}(\gamma)\times \RR
\end{cases}
\end{equation}
then $z\equiv 0$. Except for the example of Section~\ref{section-appli1}, the property is often difficult to obtain. The purpose of this section is to gather several results yielding this property.

The first known result has been proved by Ruiz in~\cite{Ruiz}. It stated the unique continuation property in a bounded domain $\Omega \subset\RR^d$ as soon as the support of $\gamma$ contains a neighborhood of the boundary $\partial\Omega$. This result has been generalized in~\cite{K-K} (see also~\cite{Las-Tri-Zha} for Neumann boundary conditions). However, this kind of results is not relevant in this paper. Indeed, their geometric settings implies the uniform decay of the semigroup $e^{At}$ and we are interested here in cases where it is not satisfied. We need sharper results.


\subsection{Unique continuation with coefficients analytic in time}

A very general unique continuation property holds if the coefficients of a linear wave equation as~\eqref{eq-UCP2} are analytic in time. This is a consequence of local continuation results proved by H\"ormander in~\cite{Hormander_uniq_analytic} and generalized by Tataru in~\cite{Tataru_uniq_analyticJMPA} and also independently proved by Robbiano and Zuily in~\cite{Rob-Zui}. These results concern in fact a very general setting but we restrict here the statements at the case of the wave equation. The application to the wave equation and the proof that the local results yield a global one are classical and straightforward, see for example~\cite[Corollary~3.2]{RC} for the details.

\begin{theo}\label{th-RZ}
{\emph{Robbiano--Zuily, H\"ormander (1998).}}
Let $T>0$ (or $T=+\infty$) and let $b$, $c$ and $d$ be smooth coefficients. Assume moreover that $b$, $c$ and $d$ are analytic in time and that $z$ is a strong solution of
\begin{equation}\label{eq-RZ}
\begin{cases}
\partial^2_{tt} z &=\Delta z + b(x,t)\partial_t z+c(x,t).\grad z+d(x,t)z\quad(x,t)\in\Omega\times (-T,T)\\
z_{|\partial\Omega}(x,t)&= 0\mkern298mu(x,t)\in\partial\Omega\times \RR\,.
\end{cases}
\end{equation}
Let $\Oc$ be a non-empty open subset of $\Omega$ and assume that $z(x,t)=0$ in $\Oc\times (-T,T)$. Then $z(x,0)\equiv 0$ in $\Oc_T=\{x_0\in\Omega\,,\,d(x_0,\Oc)<T\}$.

As consequences if $z\equiv 0$ in $\Oc\times (-T,T)$ and $\overline
\Oc_T=\Omega$, then $z\equiv 0$ everywhere.
\end{theo}



\subsection{Unique continuation through pseudo-convex surfaces without boundary}\label{subsectUCPwithout}
If the coefficients of~\eqref{eq-RZ} are not analytic in time, the geometry of the problem is more constrained. However, it could still include cases where the geometric control condition of~\cite{BLR} does not hold and thus where the semigroup $e^{At}$ is not uniformly stable, see the examples below.

We consider here H\"ormander framework (see~\cite{Hormander} for example). The principal symbol of the differential operator of~\eqref{eq-UCP2} is of order two and writes locally
\[
p(x,t,\xi,\tau)=\xi\trans.A(x).\xi-|\tau|^2
\]
where $A(x)$ is a smooth family of positive definite symmetric matrices coding the Beltrami Laplacian operator in a local chart. Let $\psi(x,t,\xi,\tau)$ be a locally $\Cc^1-$function, we introduce the Poisson bracket
\[
H_p(\psi)=\{p,\psi\}= \grad_\xi p \grad_x \psi + \partial_\tau p \partial_t
\psi - \grad_x p \grad_\xi \psi - \partial_t p \partial_\tau \psi\,.
\]
Let $\psi$ be a smooth function defined in a neighborhood $\Oc\subset
\RR^{d+1}$ of $(x_0,t_0)$. Assume that $(\grad_x \psi,\partial_t
\psi)(x_0,t_0)\neq 0$ so that $\Sigma=\{(x,t),\psi(x,t)=0\}$ defines a smooth hypersurface near $(x_0,t_0)$.

\begin{defi}\label{defi-pseudoconvex}
The local hypersurface $\Sigma$ is said to be \emph{non-characteristic} at $(x_0,t_0)$ if
\[
p(x_0,t_0,\grad \psi(x_0,t_0),\partial_t \psi(x_0,t_0))\neq 0\,.
\]
Moreover, $\Sigma$ is said to be \emph{strongly pseudo-convex} at $(x_0,t_0)$ if for any $(\xi,\tau)\neq 0$ such that $p(x_0,t_0,\xi,\tau)=0$ and $H_p(\psi)(x_0,t_0)=0$, we have
\[
H_p^2(\psi)(x_0,t_0)>0\,.
\]
\end{defi}

Notice that the above Definition~\ref{defi-pseudoconvex} of strongly pseudo-convexity is adapted to the case of a real differential operator of order two. Thus it is perfectly adapted to the situation of this paper where the wave operator is $p=\xi\trans.A(x).\xi-|\tau|^2$. However, we emphasize that, in the general case, the assumption of pseudo-convexity is more complex, see~\cite{Hormander}.

The geometrical interpretation of Definition~\ref{defi-pseudoconvex} is as follows. First, the fact that the surface is non-characteristic says that $|\partial_t \psi|^2 \neq \grad\psi\trans.A(x).\grad \psi$. This means that the surface is not moving at the exact same speed as the sound waves.

The pseudo-convexity is slightly more involved. Consider the total Hamiltonian flow $\sigma\mapsto\varphi_\sigma$ defined by
\[
\varphi_0(x,t,\xi,\tau)=(x,t,\xi,\tau)\;\partial_\sigma
\varphi_\sigma(x,t,\xi,\tau)=(\grad_\xi p,\partial_\tau
p,-\grad_x p,-\partial_t p)(\varphi_\sigma)\,.
\]
Since $p$ is independent of $t$, $\tau$ is constant and thus $t(\sigma)=t-2\tau\sigma$, meaning that $\sigma$ is a simple new parametrization of time. Moreover, $\grad_\xi p$ and $\grad_x p$ are independent of $t$ and $\tau$. Thus, $(x,\xi)(\sigma)$ follows the geodesic flow
\[
\partial_\sigma (x,\xi)=(\grad_\xi g^*,-\grad_x g^*)(x,\xi)
\]
where $g^{*}(x,\xi)=\xi\trans.A(x).\xi$ is the symbol of the local cometric. Assume $p(x,t,\xi,\tau)=0$ at $\sigma=0$. The Hamiltonian being conserved, we always have $p(x,t,\xi,\tau)=0$ and $|\tau|^2=\xi\trans A(x)\xi$ is constant in $\sigma$: the point $x(\sigma)$ is moving along a geodesic of the metric at a speed which is of constant norm $|\tau|$ with respect to the metric. Let $h$ be a function of $(x,t,\xi,\tau)$, then the Poisson bracket $\{p,h\}$ is the derivative at $\sigma=0$ of $h(\varphi_\sigma(x,t,\xi,\tau))$. Thus
\[
\{p,h\}=\{g^*,h\} -2\tau \partial_t h=\{g^*,h\} + \partial_\sigma h
\]
where $\{g^*,h\}$ is the derivative along the geodesic $(x,\xi)(\sigma)$ of the metric starting at $x$ with speed $\xi$. Thus, the strongly pseudo-convexity condition $H_p\psi=0\Rightarrow H_p^2 \psi>0$ means that if a geodesic of the surface is tangent to $\Sigma$ in the space-time sense, then it must be contained in a non-degenerated sense in the half-space $\psi(x,t)>0$ for $t\neq0$. Finally notice that if $\psi$ does not depend on time $t$ (as in Definitions~\ref{defi-foliation1} and~\ref{defi-foliation2}), then the strongly pseudo-convexity is a classical strong convexity: if a classical geodesic of the metric $g$ is tangent to the surface $\psi(x)=0$, it must be contained in the half-space $\psi(x)>0$.

\cite[Theorem~28.4.3]{Hormander} mainly comes from~\cite{LR} and is stated as follows.

\begin{theo}\label{th-hormander}
\emph{Lerner and Robbiano (1985), H\"ormander.}
Let $\Oc$ be a small open neighborhood of a point $(x_0,t_0)$ in $\RR^d\times
\RR$ and let $A(x)$ be a smooth family of positive definite symmetric matrices. Let $b$, $c$ and $d$ be bounded coefficients. Assume that $z$ is a mild solution of
\begin{equation}\label{eq-LR}
\partial^2_{tt} z=\div A(x)\grad z + b(x,t)\partial_t z+c(x,t).\grad z+d(x,t)z\;(x,t)\in\Oc\,.
\end{equation}
Let $\Sigma=\{(x,t),\psi(x,t)=0\}$ be a smooth surface containing $(x_0,t_0)$ which is non-characteristic and strongly pseudo-convex in the sense of Definition~\ref{defi-pseudoconvex}.

Then, if $u(x,t)=0$ for all $(x,t)\in\Oc$ such that $\psi(x,t)\geq 0$, we have $u(x,t)\equiv 0$ in a neighborhood of $(x_0,t_0)$.
\end{theo}

Theorem~\ref{th-hormander} states a local unique continuation property through pseudo-convex surfaces. To use it, it is more convenient to have a global version. This kind of global foliation has already been introduced in~\cite{Stefanov-Uhlmann} by Stefanov and Uhlmann.

\begin{defi}\label{defi-foliation1}
A family of surfaces $(\Sigma_\lambda)_{\lambda\in[0,1)}$ is an \emph{oriented pseudo-convex foliation without boundary} in a compact manifold $\Omega$ if:
\begin{enumerate}\romanenumi
\item \label{listdef6.4.1}the family of surfaces is smooth in the sense that it is locally described as level sets $\{x,\psi_\lambda(x)=0\}$ where $(x,\lambda)\mapsto\psi_\lambda(x)$ is a local smooth function with $\grad_x\psi_\lambda\neq 0$.
\item \label{listdef6.4.2}each surface is globally oriented in the sense that there exist disjoint sets $\Sigma_\lambda^\pm$ such that locally $\{x\in\Omega,\pm\psi_\lambda(x)> 0\}\subset \Sigma_\lambda^\pm$ and such that $\Omega=\Sigma_\lambda^- \cup
\Sigma_\lambda \cup \Sigma_\lambda^+$.
\item \label{listdef6.4.3}for each $\lambda$, $(x,t)\mapsto \psi_\lambda(x)$ is pseudo-convex in the sense of Definition~\ref{defi-pseudoconvex} as a function independent of $t$. Equivalently, $\Sigma_\lambda^-$ is locally strictly convex in a neighborhood of its boundary $\Sigma_\lambda$ for the metric $g$: for each $x\in\Sigma_\lambda$, a geodesic through $x$ which is tangent at $\Sigma_\lambda$ is locally included in $\Sigma_\lambda^+$, $x$ excepted.
\item \label{listdef6.4.4}the surfaces $\Sigma_\lambda$ are compact and have no boundary or equivalently do not meet $\partial\Omega$.
\end{enumerate}
\end{defi}

A typical example of such oriented pseudo-convex foliation without boundary is given in Figure~\ref{fig-foliation-1}.

\begin{figure}[ht]
\includegraphics{Figures/ex-foliation-1-legende.pdf} 
\caption{A oriented pseudo-convex foliation without boundary in the sphere $\SS^2$. The surfaces $\Sigma_\lambda$ forms a smooth family of vertical circles inside a hemisphere and $\Sigma_\lambda^+$ and $\Sigma_\lambda^-$ are respectively the large and the small spherical caps. The geodesics beings the great circles, the ones which are tangent to a surface $\Sigma_\lambda$ stay inside~$\Sigma_\lambda^+$.\label{fig-foliation-1}
}
\end{figure}

By a classical argument, we may state a global version of Theorem~\ref{th-hormander} as follows.

\begin{theo}\label{th-hormander-global}
Let $\Omega$ be a smooth compact manifold (with or without boundaries) and let $\omega\subset\Omega$ be an open set. Assume that there exists an oriented pseudo-convex foliation without boundary $(\Sigma_\lambda)_{\lambda\in[0,1)}$ of $\Omega$ in the sense of Definition~\ref{defi-foliation1}. Also assume that $\Sigma_0^+\subset \omega$ and $\omega\cup \left(\bigcup_{\lambda\in[0,1)}
\Sigma_\lambda^+\right)$ covers $\Omega$ up to a set of zero measure.

Let $b$, $c$ and $d$ be bounded coefficients. Assume that $z$ is a global mild solution of
\[
\begin{cases}
\partial^2_{tt} z &=\Delta z + b(x,t)\partial_t z+c(x,t).\grad z+d(x,t)z\quad(x,t)\in\Omega\times \RR\\
z_{|\omega}(x,t)&= 0\mkern297mu(x,t)\in\omega\times\RR
\end{cases}
\]
with any suitable boundary conditions on $\partial\Omega$ such that the wave equation is well-posed and where $\Delta$ is the Laplace--Beltrami operator related to $\Omega$.

Then $z\equiv 0$ everywhere.
\end{theo}


\begin{proof}
We will show that $z(\cdot,t=0)$ vanishes in $\omega\cup\Sigma^+_\lambda$, for all $\lambda\in [0,1)$, which shows that $z(t=0)\equiv 0$ and thus that $z\equiv 0$ due to the uniqueness properties of the wave equation. By assumption, $z\equiv 0$ in $\Sigma_0^+\subset \omega$. Let $\lambda_0\in (0,1)$ and let $h_{\alpha,T}(t)=\alpha(1-t^2/T^2)$. We consider the family of surfaces $t\in [-T,T]\mapsto \Sigma_{h_{\alpha,T}(t)}$ which is locally parametrized by functions $(x,t)\mapsto \psi_{h_{\alpha,T}(t)}(x)$. Notice that it is a smooth family of smooth surfaces since due to Assumption~\eqref{listdef6.4.1} of Definition~\ref{defi-foliation1}. Also notice that the larger is $T$, the smaller are the derivatives of these functions with respect to $t$. By assumption, each function $x\mapsto \psi_\lambda(x)$ is non-characteristic and strongly pseudo-convex as a function independent of $t$. By compactness, there exists $T$ large enough such that $(x,t)\mapsto \psi_{h_{\alpha,T}(t)}(x)$ defines local surfaces which are non-characteristic and pseudo-convex for all $\alpha\in[0,\lambda_0]$ and $t\in[-T,T]$. The parameter $T$ is fixed in the remaining part of the proof and we may omit it in the notations.

Notice that, for any $\alpha$, the family of set $t\in[-T,T]\mapsto
\Sigma^+_{h_\alpha(t)}$ starts inside $\omega$ at $t=-T$ and finishes inside $\omega$ at $t=T$. Moreover, for any small $\alpha$, these sets always stay inside $\omega$ where $z$ vanishes. Assume that there exist $(x,t)$ and $\alpha\in[0,\lambda_0]$ such that $x\in
\Sigma^+_{h_\alpha(t)}$ and $z(x,t)\neq 0$. We set
\begin{equation}\label{eq-demo-global}
\alpha_0=\min\left\{\alpha\in (0,\lambda_0]\,,\,
\exists t\in[-T,T],\exists x\in\Sigma^+_{h_{\alpha}(t)}\text{ such that }z(x,t)\neq 0\right\}\,.
\end{equation}
By continuity, we know that $z(x,t)=0$ for all $x\in
\Sigma^+_{h_{\alpha_0}(t)}$. Moreover, there exists $t_0\in (-T,T)$ and $x_0\in\Sigma_{h_{\alpha_0}(t_0)}$ such that $z$ is not identically zero in any neighborhood of $(x_0,t_0)$. Indeed, otherwise, by compactness, we may extend the set where $z$ vanishes and contradict~\eqref{eq-demo-global}.

To conclude, it remains to use the local unique continuation property of Theorem~\ref{th-hormander} at $(x_0,t_0)$ with the time-space surface defined by $(x,t)\mapsto \psi_{h_{\alpha_0}(t)}(x)$. The continuation implies that $z$ vanishes near $(x_0,t_0)$ which contradicts the construction. Thus, $z(x,t)=0$ for all $t\in [-T,T]$ and $x\in \bigcup_\alpha \Sigma^+_{h_\alpha(t)}$. In particular $z(\cdot,\,t=0)\equiv 0$ in $\bigcup_{\alpha}\Sigma^+_{h_\alpha(0)}=\bigcup_{\lambda\leq\lambda_0}
\Sigma^+_\lambda$. Since these arguments hold for all $\lambda_0<1$ and since $\omega\cup_{\lambda\in [0,1)} \Sigma^+_\lambda$ is $\Omega$ up to a set of measure zero, we have that $z(\cdot,\,t=0)\equiv 0$ in $\Omega$. Well-posedness of the linear wave equation concludes that $z\equiv 0$ everywhere.
\end{proof}


Notice that, as it is stated, this unique continuation result needs an infinite time to be efficient, where Theorem~\ref{th-RZ} only need a finite explicit time. In fact, a careful look to the proof shows that a finite time is sufficient once we know that the family of surface is pseudo-convex and non-characteristic in a uniform way. However, such a bound of convexity is difficult to obtain in general cases and may be even impossible as for the example studied in Section~\ref{section-appli2}.


A typical example of application is given in Figure~\ref{fig-foliation-1}: if $\Omega$ is a sphere and $\omega$ covers more than an hemisphere, then if $z$ is a global solution of a linear wave equation which vanishes in $\omega$ for all times, then $z\equiv 0$. Notice that, in this case, the family of surfaces is uniformly pseudo-convex and the unique continuation holds in fact in finite time even if $z$ is not a global in time solution.



